% 原书物理页=印刷页2--4；1.1 全节（1.1.1--1.1.10；其中 1.1.7 证明的结尾与 1.1.8--1.1.10 位于第 4 页，1.2 节自第 4 页中部开始）.
\chapter{Lie 群}\label{chap:lie-1}

在本章中，我们使用\newterm{lie-differential-geometry}{微分几何}{differential geometry}的基本词汇与概念. 所有流形及流形间的态射都将是 $C^{\infty}$ 类的，而且我们将把``$C^{\infty}$ 类的''简单地称作\newterm{lie-smooth}{光滑}{smooth}.

\section{定义}\label{sec:lie-1-1}

\begin{definition}\label{def:lie-1-1-1}
一个\newterm{lie-lie-group}{Lie 群}{Lie group}是一个集合 $G$，其上带有两个相容的结构：
\begin{itemize}[label=--]
\item 一个群结构，由一个合成运算 $m:G\times G\longrightarrow G,\ (x,y)\mapsto xy$ 给出；并把映射 $i:G\longrightarrow G,\ x\mapsto x^{-1}$ 称为``求逆映射''.
\item 一个\newterm{lie-smooth-manifold}{光滑流形}{smooth manifold}结构（由 $C^{\infty}$ 类图册的等价类给出，但我们无需为图册引入任何记号）.
\item 相容性：映射 $m$ 与 $i$ 都是光滑的.
\end{itemize}
\end{definition}

\begin{remark*}
$m$ 的 $C^{\infty}$ 性质并不蕴含 $i$ 的 $C^{\infty}$ 性质. 我们可以找到一些扭曲的反例.
\end{remark*}

\noindent\textbf{首几个例.}
\begin{itemize}[label=--]
\item 任何离散群（比如说可数的）.
\item 加法群 $(\mathbb{R},+)$.
\item 乘法群 $(\mathbb{R}^{\times},\cdot)$、其复类比 $(\mathbb{C}^{\times},\cdot)$，以及单位圆周 $(S^{1},\cdot)$.
\item 一般线性群
\end{itemize}

\begin{numpara}[一般线性群]\label{np:lie-1-1-2}
它们在某种意义上是 Lie 群的基本例子. 我们将以两种形式遇到它们.

\emph{矩阵形式}. $G=\mathrm{GL}_n(\mathbb{R})$，赋予它作为 $\mathrm{M}_n(\mathbb{R})\simeq\mathbb{R}^{n^{2}}$ 中开集的 $C^{\infty}$ 结构. 乘法 $(A,B)\mapsto AB$ 是 $A$ 与 $B$ 的元素的多项式，因而是 $C^{\infty}$ 的. 求逆映射 $A\mapsto A^{-1}=\frac{1}{\det A}\operatorname{com}(A)$ 是解析的，从而尤其是 $C^{\infty}$ 的.

\emph{内蕴形式}. 若 $V$ 是一个有限维 $\mathbb{R}$-向量空间，则令 $\mathrm{GL}(V):=\operatorname{Aut}_{\mathbb{R}}(V)$. 当然，若 $n$ 是 $V$ 的维数，则 $V$ 的一个基的选取给出一个双射 $\mathrm{GL}(V)\xrightarrow{\ \sim\ }\mathrm{GL}_n(\mathbb{R})$，按下面这个定义的意义，它是 Lie 群的一个同构.

同样，$\mathrm{GL}_n(\mathbb{C})$ 自然地是一个 Lie 群，并且若 $V$ 是一个 $\mathbb{C}$-向量空间，则 $\mathrm{GL}(V):=\operatorname{Aut}_{\mathbb{C}}(V)$ 是一个 Lie 群.
\end{numpara}

\begin{remark}\label{rem:lie-1-1-3}
设 $G$ 是一个 Lie 群，$g\in G$. 我们把 $L_{g}:G\longrightarrow G,\ x\mapsto gx$ 称为由 $g$ 给出的``左平移''. 它是 $G$ 的一个\newterm{lie-diffeomorphism}{微分同胚}{diffeomorphism}，因而诱导出各\newterm{lie-tangent-space}{切空间}{tangent space}之间的一个同构 $d_{e}L_{g}:T_{e}G\xrightarrow{\ \sim\ }T_{g}G$. 特别地，$G$ 的所有\newterm{lie-connected-component}{连通分支}{connected component}都具有相同的维数.
\end{remark}

\begin{definition}\label{def:lie-1-1-4}
Lie 群之间的一个\newterm{lie-lie-group-morphism}{态射}{morphism of Lie groups}（有时我们简单地称之为光滑态射）是一个映射 $G\xrightarrow{\ \varphi\ }G'$，它光滑且与群运算交换. 这样的态射，若它作为 $C^{\infty}$ 映射是\newterm{lie-embedding}{嵌入}{embedding}、\newterm{lie-immersion}{浸入}{immersion}或\newterm{lie-submersion}{淹没}{submersion}，就相应地称它为嵌入、浸入或淹没. 若它具有逆映射，则称它是一个\newterm{lie-isomorphism}{同构}{isomorphism}.
\end{definition}

\noindent\textbf{例.}
\begin{itemize}[label=--]
\item 行列式 $\det:\mathrm{GL}_n(\mathbb{R})\longrightarrow\mathbb{R}^{\times}$.
\item 指数映射 $(\mathbb{R},+)\xrightarrow{\ \exp\ }(\mathbb{R}_{>0}^{\times},\cdot)$ 是一个同构，其逆为对数.
\item 直线在圆周上的缠绕：$\theta\in(\mathbb{R},+)\mapsto\exp(2i\pi\theta)\in S^{1}$ 是一个满的\newterm{lie-local-diffeomorphism}{局部微分同胚}{local diffeomorphism}.
\item 直线在\newterm{lie-torus}{环面}{torus}上的缠绕：$\theta\in(\mathbb{R},+)\mapsto(\exp(2i\pi\theta),\exp(2i\pi\theta\alpha))\in S^{1}\times S^{1}$. 若 $\alpha\in\mathbb{R}\setminus\mathbb{Q}$，则这是一个单的浸入，但不是嵌入，因为其像是稠密的.
\end{itemize}

\noindent\textbf{值得注意的特殊情形.}
\begin{itemize}[label=--]
\item 态射 $\lambda:(\mathbb{R},+)\longrightarrow G$ 被称为（不严格地）\newterm{lie-one-parameter-subgroup}{单参数子群}{one-parameter subgroup}.
\item 形如 $\rho:G\longrightarrow\mathrm{GL}(V)$ 的态射称为 $G$ 在 $V$ 上的（线性）\newterm{lie-representation}{表示}{representation}. 空间 $V$ 也称为表示 $\rho$ 的``底空间''.
\end{itemize}

\begin{remark}\label{rem:lie-1-1-5}
若 $\varphi$ 是一个 Lie 群态射，则公式 $d_{h}\varphi=d_{e'}L_{\varphi(h)}\circ d_{e}\varphi\circ(d_{e}L_{h})^{-1}$ 表明 $\varphi$ 具有\newterm{lie-constant-rank}{常秩}{constant rank}.
\end{remark}

\begin{definition}\label{def:lie-1-1-6}
Lie 群 $G$ 的一个\newterm{lie-lie-subgroup}{Lie 子群}{Lie subgroup}是指一个既是\newterm{lie-submanifold}{子流形}{submanifold}又是子群的子集 $H$. 等价地，含入映射 $H\hookrightarrow G$ 是一个 Lie 群的嵌入.
\end{definition}

\noindent\textbf{例.}
\begin{itemize}[label=--]
\item $H=\mathbb{Z}$（在 $G=\mathbb{R}$ 中）.
\item $H=S^{1}$（在 $G=\mathbb{C}^{\times}$ 中）.
\end{itemize}

\noindent\textbf{注意！}不要把 Lie 子群与\newterm{lie-immersed-lie-group}{浸入 Lie 群}{immersed Lie group}（即 Lie 群的浸入 $i:H\hookrightarrow G$ 的像）相混淆. 例如，上面由无理数 $\alpha$ 给出的直线在环面上的缠绕提供了 $S^{1}\times S^{1}$ 的一个浸入 Lie 群，它同构于 $(\mathbb{R},+)$，但不是一个 Lie 子群.

\begin{proposition}\label{prop:lie-1-1-7}
Lie 群 $G$ 的任何 Lie 子群 $H$ 都是闭的.
\end{proposition}

\begin{proof}
这由下面两点推出.
\begin{enumerate}[label=\roman*)]
\item 由于 $H$ 是一个子流形，故 $H$ 是局部闭的（即 $H$ 的每一点 $h$ 在 $G$ 中都有一个邻域 $\mathcal{U}$，使得 $H\cap\mathcal{U}$ 在 $\mathcal{U}$ 中是闭的）.
\item 更一般地，设 $G$ 是一个\newterm{lie-topological-group}{拓扑群}{topological group}，$H$ 是一个局部闭子群. 则 $H$ 是闭的. 事实上，设 $g$ 属于 $H$ 在 $G$ 中的闭包 $\overline{H}$：
\begin{itemize}[label=--]
\item 选取 $G$ 中 $e$ 的一个开邻域 $\mathcal{U}$，使得 $\mathcal{U}\cap H$ 在 $\mathcal{U}$ 中是闭的. 必要时用 $\mathcal{U}\cap\mathcal{U}^{-1}$ 代替 $\mathcal{U}$，可设 $\mathcal{U}=\mathcal{U}^{-1}$（此时称 $\mathcal{U}$ 是\emph{对称的}）.
\item $g\mathcal{U}$ 是 $g$ 的一个开邻域，故 $g\mathcal{U}\cap H\neq\varnothing$. 在这个交中取一个 $h$；于是 $g\in h\mathcal{U}^{-1}\cap\overline{H}=h\mathcal{U}\cap\overline{H}$. 由于 $h\mathcal{U}$ 是开集，$h\mathcal{U}\cap\overline{H}$ 就是 $h\mathcal{U}\cap H$ 在 $h\mathcal{U}$ 中的闭包.
\item 然而，由 $\mathcal{U}$ 的取法，$h\mathcal{U}\cap H=h(\mathcal{U}\cap H)$ 在 $h\mathcal{U}$ 中是闭的.
\end{itemize}
\end{enumerate}
\end{proof}

下面的定理给出一个相当惊人的逆命题，它是众多例子的来源.

\begin{theorem}[Cartan, von Neumann]\label{thm:lie-1-1-8}
Lie 群 $G$ 的任何\newterm{lie-closed-subgroup}{闭子群}{closed subgroup} $H$ 都是一个 Lie 子群. 更确切地说，在 $H$ 上存在一个 Lie 群结构（必然唯一），使得 $H$ 成为 $G$ 的一个子流形.
\end{theorem}

我们将在后面证明这个定理. 现在给出它的一个同样惊人的推论：

\begin{corollary}\label{cor:lie-1-1-9}
设 $H$ 与 $G$ 为两个 Lie 群. 任何连续的\newterm{lie-group-homomorphism}{群同态}{group homomorphism} $\varphi:H\longrightarrow G$ 都是光滑的.
\end{corollary}

\begin{proof}
考虑 $\varphi$ 的图像映射：
\[
\Gamma_{\varphi}:\quad
\begin{array}{rcl}
H & \longrightarrow & H\times G\\
h & \longmapsto & (h,\varphi(h))
\end{array}.
\]
它是 $H$ 到其像 $\Gamma_{\varphi}(H)$ 上的一个\newterm{lie-homeomorphism}{同胚}{homeomorphism}，其逆由第一投影 $H\times G\xrightarrow{\ pr_{1}\ }H$ 在 $\Gamma_{\varphi}(H)$ 上的限制 $pr_{1}|_{\Gamma_{\varphi}(H)}$ 给出. 此外，等式 $\Gamma_{\varphi}(H)=\{(h,g)\in H\times G,\ \varphi(h)=g\}$ 表明 $\Gamma_{\varphi}(H)$ 是闭的. 根据前面的定理，它因此是 $H\times G$ 的一个 Lie 子群. 而 $pr_{1}$ 显然是光滑的，故 $pr_{1}|_{\Gamma_{\varphi}(H)}$ 是一个光滑且具有常秩的同胚（见上文相关注）. 常秩定理蕴含它是一个微分同胚. 由此推出 $\Gamma_{\varphi}$ 是一个光滑态射，从而若记 $pr_{2}$ 为第二投影，则 $\varphi=pr_{2}\circ\Gamma_{\varphi}$ 也是光滑的.
\end{proof}

\begin{remark}\label{rem:lie-1-1-10}
值得强调的是，在上面的证明中用到了秩定理的如下推论：任何双射的 Lie 群态射都是一个同构. 事实上，只需证明 $\varphi$ 在点 $e$ 处的微分 $d_{e}\varphi:T_{e}H\longrightarrow T_{e}G$ 是可逆的（此时它在每一点处都可逆）. 若 $d_{e}\varphi$ 不是单射，则常秩定理立即蕴含 $\varphi$ 在 $e$ 的某个邻域上不是单射，这与 $\varphi$ 的单性矛盾. 另一方面，若 $d_{e}\varphi$ 不是满射\footnote{译注：原文此处作 ``di $d_{e}\varphi$''，疑为 ``si $d_{e}\varphi$'' 之误.}，则常秩定理蕴含 $e$ 的某个适当邻域在 $\varphi$ 下的像具有空内部. 由 Baire 定理可知，$\varphi$ 的像具有空内部，这与 $\varphi$ 的满性矛盾.
\end{remark}
